\documentclass[11pt, notitlepage, letterpaper]{article}
\usepackage{array} % conditions
\usepackage{amsmath} % dfrac
\usepackage{braket} % dirac notation
\begin{document}
\centering Depolarizing (Pauli) Error Channel
\\ \noindent For some quantum state $\rho$ and
\\ \noindent physical qubit with error probability $p$:
\begin{align}E_{\text{pauli}}(\rho) & = (1-p)\rho + \dfrac{p}{3}(X\rho X + Y\rho Y + Z\rho Z)
\end{align}


\centering
\noindent For some quantum state $\rho$ and
\\ \noindent physical qubit with error probability $p$:
\begin{align}E_{\text{pauli}}(\rho) & = (1-p)\rho + \dfrac{p}{3}(X\rho X + Y\rho Y + Z\rho Z)
\end{align}

An erasure occurs when a qubit no longer detectable by the measurement device, so the quantum state effectively ``lost'' or ``erased''. If we extend the mathematical representation of a qubit Hilbert space with some $\ket{e}$ that is orthogonal to the qubit state $\{\ket{0}, \ket{1}\}$, then we can define the erasure channel as:
\begin{align}E_{\text{erase}}(\rho_j) & = (1-p)\rho_j + p\ket{e}\bra{e}
\end{align}
% Sample uniformly at random with probability:
% \\ $1-p$ no error
% \\ $p/3$ error in $\set{X, Y, Z}$


\noindent qubit at location $j$ leaves the computational subspace \\ flagged as
``erased'', we model with orthogonal state $\ket{e}$:
\[
    E_{\mathrm{erase}}(\rho_j) = (1-p)\rho_j + p\,\ket{e}\!\bra{e}.
\]
qubit $j$ reset to maximally mixed state:
$\rho_j = \tfrac{1}{2}\,I$.

\end{document}